\documentclass[11pt,a4paper]{article}
\usepackage[margin=2.5cm]{geometry}
\usepackage{amsmath,amssymb}
\usepackage{booktabs}
\usepackage{longtable}
\usepackage{hyperref}
\usepackage{xcolor}
\usepackage{microtype}
\usepackage{parskip}
\usepackage{listings}
\usepackage{graphicx}

\hypersetup{colorlinks=true, linkcolor=blue, citecolor=blue, urlcolor=blue}

\lstset{basicstyle=\ttfamily\small, breaklines=true, frame=single,
        backgroundcolor=\color{gray!8}}

\title{\textbf{Reproducing ESM-2 650M Masked Marginal Scoring on ProteinGym}\\[0.4em]
\large A reproducibility study with empirical timing analysis}
\author{Reproducibility run: \texttt{esm2-inference-opt}\\
\small Hardware: NVIDIA A100-SXM4-40GB \quad
Modal workspace: \texttt{tyronita}\\
\small fair-esm SHA: \texttt{2b369911} \quad
transformers: \texttt{4.44.0} \quad
consistency gate: $\rho = +0.999973$}
\date{September 2026}

\begin{document}
\maketitle
\tableofcontents
\vspace{1em}
\noindent\textbf{Reproducibility statement.}
All code, pinned dependency SHAs, and results are in
\url{https://github.com/Tyronita/esm2-inference-opt}.
The reference implementation is \texttt{refs/fair-esm/} at commit
\texttt{2b369911bb5b4b0dda914521b9475cad1656b2ac} (fair-esm v2.0.1).
Benchmark data: ProteinGym v1.3, \texttt{refs/ProteinGym/} at commit
\texttt{144fe22b07dfaeec2b366f2346203a9838a55b4c}.

% ─────────────────────────────────────────────────────────────────────────────
\section{Method: Masked Marginal Scoring}

Masked marginal scoring~\cite{meier2021language} quantifies the effect of a
mutation by comparing how likely the mutant residue is versus the wildtype
residue, \emph{given every other position in the sequence as context}.

For a protein sequence $\mathbf{x} = (x_1, \ldots, x_L)$ and a set of
mutations $\mathcal{M} = \{(i, x^*_i)\}$ (position, mutant residue):

\begin{equation}
  \text{Score}(\mathbf{x}, \mathcal{M}) =
    \sum_{i \in \mathcal{M}}
    \Bigl[
      \log p\!\left(x^*_i \,\big|\, \mathbf{x}_{-i}\right)
      -
      \log p\!\left(x_i \,\big|\, \mathbf{x}_{-i}\right)
    \Bigr]
  \label{eq:score}
\end{equation}

where $\mathbf{x}_{-i}$ denotes the wildtype sequence with position $i$
replaced by the special \texttt{[MASK]} token. Each term is a
\emph{log-likelihood ratio}: positive means the model prefers the mutant;
negative means the model prefers the wildtype.

\paragraph{Computational cost.}
Evaluating Equation~\ref{eq:score} for all variants in an assay requires
$|\{i : i \in \text{some variant}\}|$ forward passes — one mask per unique
mutated position. For a typical assay of length $L$, this is $L$ passes,
each processing a sequence of $L$ tokens. The total cost scales as
$\mathcal{O}(L^2)$ in the number of residues.

\paragraph{Implementation.}
Our implementation (\texttt{scoring/masked\_marginals.py}) uses
\texttt{transformers==4.44.0} with \texttt{torch\_dtype=float16} on CUDA.
We verified it reproduces the reference fair-esm (fp32) scores with
Spearman $\rho = +0.999973$ on 200 SNCA single mutants
(consistency check: \texttt{check\_consistency.py}).

% ─────────────────────────────────────────────────────────────────────────────
\section{ESM-2 Architecture}

ESM-2~\cite{lin2023esm2} is a family of transformer-based protein language
models trained on UniRef50 (250M sequences) using masked language modelling.
Key architectural choices (from \texttt{refs/fair-esm/esm/model/esm2.py}):

\begin{itemize}
  \item \textbf{Rotary positional embeddings (RoPE)} — no absolute positional
        embeddings; relative position encoded in attention.
  \item \textbf{Pre-LayerNorm} — normalisation before each sub-layer.
  \item \textbf{GELU activations} in the feed-forward network.
  \item \textbf{Token dropout} during training (masked at input level).
  \item \textbf{Vocabulary size} = 33 (20 standard AAs + special tokens).
  \item \textbf{Maximum context} = 1024 tokens (1022 residues + BOS + EOS).
        Sequences longer than 1022 AA require windowing.
  \item \textbf{FFN dimension} = $4 \times d_\text{model}$.
\end{itemize}

\begin{table}[h]
\centering
\caption{ESM-2 model family. ProteinGym $\rho$ from
\texttt{refs/ProteinGym/benchmarks/.../Summary\_performance\ldots.csv}.
\textbf{Bold}: model used in this study.}
\label{tab:models}
\begin{tabular}{lrrrrrl}
\toprule
Model ID & Layers & $d$ & Heads & FFN & Params & PG $\rho$ \\
\midrule
esm2\_t6\_8M\_UR50D    &  6 &  320 & 20 &  1280 &   8M & 0.226 \\
esm2\_t12\_35M\_UR50D  & 12 &  480 & 20 &  1920 &  35M & 0.321 \\
esm2\_t30\_150M\_UR50D & 30 &  640 & 20 &  2560 & 150M & 0.387 \\
\textbf{esm2\_t33\_650M\_UR50D} & \textbf{33} & \textbf{1280} & \textbf{20} & \textbf{5120} & \textbf{650M} & \textbf{0.414} \\
esm2\_t36\_3B\_UR50D   & 36 & 2560 & 40 & 10240 &   3B & 0.406 \\
esm2\_t48\_15B\_UR50D  & 48 & 5120 & 40 & 20480 &  15B & 0.400 \\
\bottomrule
\end{tabular}
\end{table}

\noindent\textbf{Observation.} The 650M model achieves the highest ProteinGym
score despite being smaller than 3B and 15B variants. Larger models may
overfit to evolutionary co-occurrence patterns that do not generalise to
experimental fitness measurements.

% ─────────────────────────────────────────────────────────────────────────────
\section{Benchmark: ProteinGym}

ProteinGym~\cite{notin2023proteingym} is a large-scale benchmark of
deep mutational scanning (DMS) assays measuring protein fitness.
We use the \textbf{substitution benchmark}: 217 assays, 2,465,767 total
variants, covering diverse proteins, organisms, and fitness phenotypes.

\paragraph{Metric.}
For each assay, we compute the \textbf{Spearman rank correlation}
$\rho_s$ between the predicted scores (Equation~\ref{eq:score}) and the
experimentally measured fitness values. Spearman $\rho$ is chosen because
fitness scales vary wildly across assays; rank correlation is invariant to
monotone rescaling. The aggregate is a mean over 217 per-assay $\rho_s$.

\paragraph{Published baseline.}
ESM-2 650M with masked marginals scores $\rho = 0.414 \pm 0.012$ (mean
$\pm$ bootstrap SE) across 217 assays~\cite{notin2023proteingym}.
This is the target we reproduce.

% ─────────────────────────────────────────────────────────────────────────────
\section{Related Work}

The following papers provide independent evidence that masked language
model log-likelihoods are predictive of protein variant fitness.

\begin{longtable}{p{3.5cm}p{9cm}}
\caption{Papers supporting the masked LM approach to variant effect prediction.}
\label{tab:related} \\
\toprule
Paper & Agreement with this work \\
\midrule
\endfirsthead
\toprule
Paper & Agreement with this work \\
\midrule
\endhead
\midrule
\multicolumn{2}{r}{\small\itshape continued\ldots}
\endfoot
\bottomrule
\endlastfoot
Rives et al.\ (2021) PNAS~\cite{rives2021biological}
  & ESM-1b (predecessor to ESM-2) shows that large protein LMs trained on
    evolutionary sequences learn representations that capture structural and
    functional properties; validates the pretraining approach at 650M scale. \\[4pt]
Elnaggar et al.\ (2021) IEEE TPAMI~\cite{elnaggar2021prottrans}
  & ProtTrans (independent T5/BERT models on UniRef) independently confirms
    that masked language model likelihoods correlate with variant effects
    across multiple benchmark datasets. \\[4pt]
Notin et al.\ (2022) ICML~\cite{notin2022tranception}
  & Tranception introduces autoregressive scoring on the same DMS assays
    used in ProteinGym; establishes the benchmark structure we target and
    shows masked LMs are competitive baselines. \\[4pt]
Kantroo et al.\ (2024) PRX Life~\cite{kantroo2024ofs}
  & One Fell Swoop (OFS) achieves near-masked-marginal accuracy with a
    single forward pass by training MLP heads on ESM-2 hidden states;
    confirms that ESM-2 representations carry sufficient signal to predict
    fitness in $\mathcal{O}(1)$ cost. \\
\end{longtable}

% ─────────────────────────────────────────────────────────────────────────────
\section{Empirical Timing Model}

The total wall-clock time for masked marginal scoring on a single assay
follows two regimes depending on sequence length $L$:

\begin{equation}
  t(L) \approx
  \begin{cases}
    \beta \cdot L^2 & L \leq 1022 \\
    \beta' \cdot L  & L > 1022 \quad\text{(windowing: each pass costs } \mathcal{O}(1022) \text{ tokens)}
  \end{cases}
  \label{eq:timing}
\end{equation}

\begin{table}[h]
\centering
\caption{Measured $\beta$ (s per AA$^2$) on different hardware.}
\label{tab:timing}
\begin{tabular}{llrl}
\toprule
Hardware & Measured on & $\beta$ (s/AA$^2$) & Note \\
\midrule
Apple M3 MPS  & SNCA L=140 (16.5s), LRRK2 L=2527 (85.8min) & $8.25 \times 10^{-4}$ & L$^2$ scaling confirmed \\
A100-SXM4-40GB & SYUA L=140 (3.7s), PRKN L=465 (12.8s), BRCA1 L=1863 (12.1s) & $3.14 \times 10^{-5}$ & median across 4 assays \\
\midrule
Speedup (A100 vs M3) & \multicolumn{3}{l}{$\approx 26\times$ per unit $L^2$} \\
\bottomrule
\end{tabular}
\end{table}

\noindent Using $\beta_\text{A100} = 3.14 \times 10^{-5}$, the predicted
cost for the full 217-assay run is:
$\sum_i L_i^2 \times 3.14 \times 10^{-5} \approx 2{,}783$\,s $\approx$ 46\,min
for the $L \leq 1022$ assays; assays with $L > 1022$ are faster due to windowing.

% ─────────────────────────────────────────────────────────────────────────────
\section{Results}
\label{sec:results}

Results are appended to this section automatically by \texttt{paper/generate\_results.py}
after each run. The table below is append-only; each row is one completed assay.

% Auto-generated by paper/generate_results.py — do not edit manually.
% Each row is one completed assay run. New rows are appended by generate_results.py.

\begin{table}[h]
\centering
\caption{Assay-level results. $\rho$: Spearman rank correlation vs.\ experiment.
$\beta$: wall\_s / L$^2$ (s/AA$^2$). Published target: $\rho = 0.414 \pm 0.012$.}
\label{tab:results}
\begin{tabular}{llrrrrl}
\toprule
Run date & Assay & $L$ & $N_\text{var}$ & $\rho$ & $\beta$ (s/AA$^2$) & Hardware \\
\midrule
% ── pd_proteins ──
2026-09-19 & SYUA_HUMAN_Newberry_2020 & 140 & 2,497 & $+0.128$ & $3.05e-5$ & NVIDIA A100-SXM4-40GB \\
2026-09-19 & SYUA_HUMAN_Newberry_2020 & 140 & 2,497 & $+0.128$ & $2.14e-5$ & NVIDIA A100-SXM4-40GB \\
2026-09-19 & PRKN_HUMAN_Clausen_2023 & 465 & 8,756 & $+0.503$ & $1.63e-6$ & NVIDIA A100-SXM4-40GB \\
2026-09-19 & BRCA1_HUMAN_Findlay_2018 & 1863 & 1,837 & $+0.514$ & $1.40e-7$ & NVIDIA A100-SXM4-40GB \\
2026-09-19 & A0A140D2T1_ZIKV_Sourisseau_2019 & 3423 & 9,576 & $+0.209$ & $9.00e-8$ & NVIDIA A100-SXM4-40GB \\
2026-09-19 & SYUA_HUMAN_Newberry_2020 & 140 & 2,497 & $+0.137$ & $1.90e-4$ & NVIDIA A100-SXM4-40GB \\
2026-09-19 & PRKN_HUMAN_Clausen_2023 & 465 & 8,756 & $+0.501$ & $5.94e-5$ & NVIDIA A100-SXM4-40GB \\
2026-09-19 & BRCA1_HUMAN_Findlay_2018 & 1863 & 1,837 & $+0.515$ & $3.48e-6$ & NVIDIA A100-SXM4-40GB \\
2026-09-19 & A0A140D2T1_ZIKV_Sourisseau_2019 & 3423 & 9,576 & $+0.216$ & $1.61e-6$ & NVIDIA A100-SXM4-40GB \\
2026-09-19 & SYUA_HUMAN_Newberry_2020 & 140 & 2,497 & $+0.137$ & $3.35e-3$ & NVIDIA A100-SXM4-40GB \\
2026-09-19 & PRKN_HUMAN_Clausen_2023 & 465 & 8,756 & $+0.501$ & $1.09e-3$ & NVIDIA A100-SXM4-40GB \\
2026-09-19 & BRCA1_HUMAN_Findlay_2018 & 1863 & 1,837 & $+0.515$ & $2.03e-5$ & NVIDIA A100-SXM4-40GB \\
\bottomrule
\end{tabular}
\end{table}


\noindent\textbf{Published target:} ESM-2 650M masked marginals,
$\rho = 0.414 \pm 0.012$ (217 assays, Notin et al.\ 2023~\cite{notin2023proteingym}).

% ─────────────────────────────────────────────────────────────────────────────
\bibliographystyle{plain}
\bibliography{refs}

\end{document}
